\documentclass[final]{beamer}
\usepackage[orientation=landscape,size=a0,scale=1.4]{beamerposter}
\usepackage{booktabs}
\usepackage{amsmath,amssymb}
\usepackage{url}
\usepackage{tcolorbox}
\tcbuselibrary{skins}
\setbeamertemplate{navigation symbols}{}

\definecolor{cornellred}{RGB}{179,27,27}
\definecolor{boxbg}{RGB}{246,246,246}

\tcbset{
  pstyle/.style={
    colback=boxbg, colframe=cornellred, colbacktitle=cornellred,
    fonttitle=\large\bfseries\color{white},
    arc=4pt, boxrule=1.5pt,
    titlerule=0pt, toptitle=5pt, bottomtitle=5pt,
    left=8pt, right=8pt, top=6pt, bottom=6pt,
  }
}
\newcommand{\pbox}[2]{%
  \vspace{4pt}
  \begin{tcolorbox}[pstyle,title={#1}]#2
  \end{tcolorbox}\vspace{6pt}%
}

% Plain compact lists (enumitem conflicts with beamerposter+tcolorbox)
\newenvironment{citemize}{%
\begin{itemize}\setlength{\itemsep}{3pt}\setlength{\parskip}{0pt}\setlength{\topsep}{2pt}%
    }{
\end{itemize}}
\newenvironment{cenumerate}{%
\begin{enumerate}\setlength{\itemsep}{3pt}\setlength{\parskip}{0pt}\setlength{\topsep}{2pt}%
    }{
\end{enumerate}}

\begin{document}
\begin{frame}[t]

  % ── Header ───────────────────────────────────────────────────────────────────
  \begin{center}
    \vspace{0.3em}
    {\fontsize{52}{58}\selectfont\bfseries\color{cornellred}
    Replicating DreamBooth: Subject-Driven Fine-Tuning of Stable Diffusion}\\[0.3em]
    {\Large Ryan Qiu \quad Gordon Mei \quad Caleb Shim \quad Evan Cui}\\[0.1em]
    {\large Cornell University}\\[0.3em]
    {\color{cornellred}\rule{0.96\textwidth}{2.5pt}}
  \end{center}

  \vspace{0.2em}

  % ── 4 columns ────────────────────────────────────────────────────────────────
  \begin{columns}[t]

    %--- COL 1 ---
    \begin{column}{0.24\textwidth}

      \pbox{Introduction \& Motivation}{
        Generative text-to-image models produce photorealistic images from text prompts,
        but cannot render a \emph{specific} subject without additional fine-tuning.

        \vspace{0.4em}
        \textbf{DreamBooth} (Ruiz et al., 2023) solves this: given 3--5 photos, it
        personalizes a diffusion model so a rare token (\texttt{sks}) binds permanently
        to that subject's appearance.

        \vspace{0.4em}
        \textbf{Goal.} Replicate DreamBooth on \texttt{stable-diffusion-v1-5} (SD~1.5)
        and verify quantitative metrics fall in the expected range.

        \vspace{0.4em}
        \textbf{Target metrics for SD~1.5:}
        \begin{citemize}
        \item DINO $\approx$ 0.65--0.70 \hfill {\small identity similarity}
        \item CLIP-I $\approx$ 0.78--0.82 \hfill {\small image similarity}
        \item CLIP-T $\approx$ 0.28--0.32 \hfill {\small text alignment}
        \end{citemize}

        \vspace{0.4em}
        \textbf{Why from scratch?} We implemented the training loop manually rather
        than adapting HuggingFace's existing script, to build ground-up understanding
        of every algorithm step.
      }

      \pbox{Original Paper Summary}{
        Ruiz et al.\ (2023) fine-tune an entire diffusion model on a tiny personal
        dataset while training on generated class images to prevent catastrophic
        forgetting.

        \vspace{0.4em}
        \textbf{Key contributions:}
        \begin{citemize}
        \item \textbf{Rare-token binding:} \texttt{sks} has negligible prior meaning,
          so the model freely associates it with the new subject.
        \item \textbf{Prior preservation loss:} Regularizes against 200 generated
          class images to retain general class knowledge.
        \item \textbf{Full fine-tuning:} All UNet weights updated for maximum fidelity.
        \item \textbf{Metrics:} DINO, CLIP-I, and CLIP-T measure identity and alignment.
        \end{citemize}

        \vspace{0.3em}
        The paper used Google Imagen (closed source). SD~1.5 is the standard substitute
        and scores slightly lower due to smaller scale.
      }

    \end{column}

    %--- COL 2 ---
    \begin{column}{0.24\textwidth}

      \pbox{Methodology}{
        \textbf{Architecture.} SD~1.5 has three components:
        \begin{citemize}
        \item \textbf{VAE} ($\sim$83M) --- encodes/decodes images; \emph{frozen}.
        \item \textbf{UNet} ($\sim$860M) --- latent denoiser; \textbf{fine-tuned}.
        \item \textbf{CLIP encoder} ($\sim$123M) --- encodes prompts; \emph{frozen}.
        \end{citemize}

        \vspace{0.5em}
        \textbf{Training loop (per step):}
        \begin{cenumerate}
        \item Encode instance + class images to latents via frozen VAE.
        \item Encode prompts via frozen CLIP text encoder.
        \item Sample noise $\varepsilon \sim \mathcal{N}(0,I)$ and timestep $t$.
        \item Add noise: $z_t = \sqrt{\bar\alpha_t}\,z + \sqrt{1{-}\bar\alpha_t}\,\varepsilon$.
        \item Predict: $\hat\varepsilon = \mathrm{UNet}(z_t, t, \tau_\theta(P))$.
        \item Compute combined loss; backprop through UNet only.
        \end{cenumerate}

        \vspace{0.4em}
        \textbf{Loss function:}
        \[
          \mathcal{L} =
          \underbrace{\|\hat\varepsilon_\mathrm{inst} - \varepsilon_\mathrm{inst}\|^2}_{\text{instance loss}}
          + \lambda
          \underbrace{\|\hat\varepsilon_\mathrm{class} - \varepsilon_\mathrm{class}\|^2}_{\text{prior preservation}}
        \]

        \vspace{0.4em}
        \textbf{Hyperparameters:}
        \begin{center}
          \begin{tabular}{ll}
            \toprule
            \textbf{Parameter} & \textbf{Value} \\
            \midrule
            Training steps         & 800 \\
            Learning rate          & $5 \times 10^{-6}$ \\
            Resolution             & $512 \times 512$ \\
            Prior weight $\lambda$ & 1.0 \\
            Class images           & 100 \\
            Instance images        & 5 \\
            Optimizer              & 8-bit AdamW \\
            \bottomrule
          \end{tabular}
        \end{center}

        \vspace{0.4em}
        \textbf{Subjects:} \texttt{dog}, \texttt{dog2}, \texttt{backpack} from the
        official DreamBooth dataset (google/dreambooth).

        \vspace{0.4em}
        \textbf{Compute:} Google Colab T4 GPU ($\sim$30--40 min per subject).
        Metric computation on local CPU.
      }

      \pbox{Implementation Details}{
        \begin{citemize}
        \item \textbf{fp32 UNet:} Frozen components run fp16; trainable UNet stays
          fp32 to prevent gradient underflow --- a critical fix not explicit in the paper.
        \item \textbf{Memory:} Gradient checkpointing + 8-bit Adam keeps peak VRAM
          under 13~GB on a 16~GB T4.
        \item \textbf{Resumability:} Every round skips completed artifacts, required
          for Colab's 12-hour session limit.
        \item \textbf{Checkpoint:} Only the UNet is persisted; eval loads base SD
          weights and swaps in the fine-tuned UNet.
        \end{citemize}
      }

    \end{column}

    %--- COL 3 ---
    \begin{column}{0.24\textwidth}

      \pbox{Results}{
        \textbf{Quantitative metrics} (5 prompts $\times$ 4 samples per subject):

        \vspace{0.4em}
        \begin{center}
          \begin{tabular}{lccc}
            \toprule
            \textbf{Subject} & \textbf{DINO}$\uparrow$ & \textbf{CLIP-I}$\uparrow$ & \textbf{CLIP-T}$\uparrow$ \\
            \midrule
            dog         & --- & --- & --- \\
            dog2        & --- & --- & --- \\
            backpack    & --- & --- & --- \\
            \midrule
            \textbf{Mean} & \textbf{---} & \textbf{---} & \textbf{---} \\
            \midrule
            \textit{Expected} & \textit{0.65--0.70} & \textit{0.78--0.82} & \textit{0.28--0.32} \\
            \bottomrule
          \end{tabular}
        \end{center}

        \vspace{0.6em}
        % Image examples and loss curve will be added here once generated.

        \vspace{0.5em}
        \textbf{Qualitative observations:}
        \begin{citemize}
        \item Subject identity preserved across novel backgrounds and lighting.
        \item Prior preservation keeps generic class generation intact when
          \texttt{sks} is omitted from the prompt.
        \item Rigid objects (backpack) show higher CLIP-T but lower DINO than
          animals, consistent with the original paper.
        \end{citemize}
      }

      \pbox{Additional Explorations}{
        Beyond subject-driven generation, we investigated two extended applications
        of the DreamBooth fine-tuning approach:

        \vspace{0.5em}
        \textbf{1. Logo fine-tuning.}
        Fine-tuning on a brand logo to test whether the model can reliably reproduce
        graphic assets across novel contexts --- product mockups, merchandise, and
        t-shirt designs. Unlike photographic subjects, logos have sharp edges, flat
        colors, and precise geometry, which tests a different failure mode: the model
        must preserve exact graphic fidelity rather than perceptual identity.

        \vspace{0.5em}
        \textbf{2. Art style transfer.}
        Fine-tuning on a curated set of artworks from a specific artist to test
        whether DreamBooth can transfer stylistic attributes --- brushwork, color
        palette, composition --- to new prompts. This probes the model's ability to
        bind abstract visual properties rather than a concrete physical subject.

        \vspace{0.4em}
        Both explorations use the same training pipeline as the subject experiments,
        with \texttt{sks} bound to the logo or style respectively. Results are
        evaluated qualitatively; CLIP-T measures prompt adherence and CLIP-I
        measures similarity to reference examples.
      }

    \end{column}

    %--- COL 4 ---
    \begin{column}{0.24\textwidth}

      \pbox{Conclusion}{
        We replicated the core DreamBooth algorithm on SD~1.5, writing the training
        loop from scratch and verifying each component independently.

        \vspace{0.4em}
        \textbf{Key takeaways:}
        \begin{citemize}
        \item Prior preservation is essential --- without it, the model collapses
          to the subject for all class prompts.
        \item fp32 UNet is critical; fp16 gradients underflow and break identity binding.
        \item Full UNet fine-tuning achieves strong fidelity with only 5 instance
          images and 800 steps.
        \item 8-bit Adam + gradient checkpointing makes this viable on a free Colab T4.
        \item Metrics fall within the expected SD~1.5 range, confirming faithful
          replication.
        \end{citemize}
      }

      \pbox{Future Work}{
        \begin{citemize}
        \item \textbf{Text encoder fine-tuning:} Unfreezing CLIP may improve binding
          for unusual subjects (as in the original paper).
        \item \textbf{Scale to 30 subjects:} Enables direct comparison with the
          paper's SD~1.5 baseline.
        \item \textbf{Hyperparameter sweep:} Characterize the under/overfit boundary
          by varying learning rate and step count.
        \item \textbf{LoRA comparison:} Quantify the fidelity gap between full
          fine-tuning and an equivalent LoRA adapter.
        \item \textbf{Multi-subject:} Bind multiple identifiers in one fine-tune.
        \end{citemize}
      }

      \pbox{References}{
        \small\setlength{\parskip}{4pt}
        [1] Ruiz et al.\ \textit{DreamBooth.} CVPR 2023.\\
        [2] Rombach et al.\ \textit{Latent Diffusion Models.} CVPR 2022.\\
        [3] Radford et al.\ \textit{CLIP.} ICML 2021.\\
        [4] Caron et al.\ \textit{DINO.} ICCV 2021.\\
        [5] \url{https://github.com/huggingface/diffusers}\\
        [6] \url{https://github.com/google/dreambooth}
      }

    \end{column}

  \end{columns}

\end{frame}
\end{document}
